% T3b: Out-of-vocabulary generalization (paper_experiments.py, analysis "oov")
% Models trained on a 90% file split; OOV = forms appearing only in the held-out 10%.
\begin{table}[t]
  \centering
  \begin{tabular}{lccc}
    \toprule
    \textbf{Representation} & \textbf{OOV forms vectorized} & \textbf{coverage}
    & \textbf{NN shares root prefix} \\
    \midrule
    FastText (char $n$-gram) & 9{,}023 / 9{,}023 & 100\% & 169 / 300 \;(56\%) \\
    word2vec (no subword)    & 0 / 9{,}023      & 0\%   & n/a \\
    \bottomrule
  \end{tabular}
  \caption{Out-of-vocabulary generalization. With models trained on a 90\% file
  split, 9{,}023 forms (23.3\% of the held-out vocabulary) are truly unseen.
  FastText synthesizes a vector for \emph{every} one from its character
  $n$-grams, and for a 300-form sample the nearest in-vocabulary neighbour shares
  the three-consonant root prefix 56\% of the time---evidence the synthesized
  vectors are morphologically coherent. \textsc{word2vec} cannot represent any
  unseen form.}
  \label{tab:oov}
\end{table}
